\usepackage[most]{tcolorbox}

%----------------------------------------------------------------------------------------
%   HIGHLIGHT ENVIRONMENT
%----------------------------------------------------------------------------------------

\newtcolorbox{highlight}{
  colback=bg,
  colframe=accent,
  coltext=fg,
  boxrule=2pt,
  boxsep=4pt,
  arc=0pt,
  outer arc=0pt
}

%----------------------------------------------------------------------------------------
%   REMARK BOX ENVIRONMENT
%----------------------------------------------------------------------------------------

\newtcolorbox{remarkbox}{
  colback=bg,
  colframe=accent,
  coltext=fg,
  boxrule=0pt,
  leftrule=2pt,
  boxsep=4pt,
  arc=0pt,
  outer arc=0pt,
}

%----------------------------------------------------------------------------------------
%   OUTLINE ENVIRONMENT
%----------------------------------------------------------------------------------------

\newtcolorbox{outline}[1][4pt]{
  enhanced,
  colback=bg,
  colframe=bg, % To essentially hide the frame, but we will draw the corners manually
  coltext=fg,
  boxrule=0pt,
  boxsep=#1,
  arc=0pt,
  outer arc=0pt,
  overlay={
    % Top left corner
    \draw[accent,line width=2pt] 
      (frame.north west) -- ++(0,-0.25*\tcbtextheight)
      (frame.north west) -- ++(0.25*\tcbtextwidth,0);
    % Bottom right corner
    \draw[accent,line width=2pt]
      (frame.south east) -- ++(0,0.25*\tcbtextheight)
      (frame.south east) -- ++(-0.25*\tcbtextwidth,0);
  }
}

%----------------------------------------------------------------------------------------
%   EDITOR COMMENT ENVIRONMENT
%----------------------------------------------------------------------------------------

\definecolor{Hisu}{HTML}{009630}
\definecolor{HisuQuestion}{HTML}{D90429}
\providecommand{\normal}{\normalsize}

\tcbset{
  hisunote/.style={
    enhanced,
    breakable,
    colback=Hisu!4!bg,
    colframe=Hisu!60,
    coltext=Hisu,
    boxrule=0.6pt,
    arc=2pt,
    outer arc=2pt,
    left=7pt, right=6pt,
    top=3pt, bottom=3pt,
    borderline west={2pt}{0pt}{Hisu},
    before skip=4pt,
    after skip=4pt,
  }
}

\newenvironment{hisublock}[1][\footnotesize]{%
  \begin{tcolorbox}[hisunote]%
  \renewcommand{\textbf}[1]{\oldtextbf{##1}}%
  #1\ignorespaces
}{%
  \end{tcolorbox}%
}

\newenvironment{hisuquestion}[1][\footnotesize]{%
  \par\begingroup\color{HisuQuestion}#1\noindent
  \renewcommand{\textbf}[1]{\oldtextbf{##1}}%
  \ignorespaces
}{%
  \par\endgroup\ignorespacesafterend
}

\newcommand{\hisu}[1]{{\small\color{Hisu}#1}}
\newcommand{\hisuq}[1]{{\small\color{HisuQuestion}#1}}

%----------------------------------------------------------------------------------------
%   THEOREM LEFT-RULE STYLE
%----------------------------------------------------------------------------------------

\tcbset{
  thmbar/.style={
    enhanced,
    breakable,
    colback=bg,
    colframe=accent,
    coltext=fg,
    boxrule=0pt,
    leftrule=2.5pt,
    toprule=0.5pt,
    bottomrule=0.5pt,
    rightrule=0.5pt,
    arc=0pt,
    outer arc=0pt,
    left=6pt,
    right=4pt,
    top=2pt,
    bottom=3pt,
    before skip=6pt,
    after skip=6pt,
  }
}

\tcolorboxenvironment{theorem}{thmbar}
\tcolorboxenvironment{definition}{thmbar}
\tcolorboxenvironment{proposition}{thmbar}
\tcolorboxenvironment{lemma}{thmbar}
\tcolorboxenvironment{corollary}{thmbar}
\tcolorboxenvironment{condition}{thmbar}
\tcolorboxenvironment{remark}{thmbar}
\tcolorboxenvironment{example}{thmbar}
\tcolorboxenvironment{question}{thmbar}

% Reset flag at end of each theorem env (in case body had no list)
\AtEndEnvironment{theorem}{\global\theoremheadnewlinefalse}
\AtEndEnvironment{definition}{\global\theoremheadnewlinefalse}
\AtEndEnvironment{proposition}{\global\theoremheadnewlinefalse}
\AtEndEnvironment{lemma}{\global\theoremheadnewlinefalse}
\AtEndEnvironment{corollary}{\global\theoremheadnewlinefalse}
\AtEndEnvironment{condition}{\global\theoremheadnewlinefalse}
\AtEndEnvironment{remark}{\global\theoremheadnewlinefalse}
\AtEndEnvironment{example}{\global\theoremheadnewlinefalse}
\AtEndEnvironment{question}{\global\theoremheadnewlinefalse}
